\section{Annotating the Information System}
\label{sec:mapping}

This section describes how the relational database is mapped to \ac{rdf} triples using \ac{r2rml}~\cite{r2rml}, producing an instance of the CineExplorer knowledge graph conforming to the ontology described in the previous section.

\subsection{Approach}

\ac{r2rml} (RDB to \ac{rdf} Mapping Language)~\cite{r2rml} is used as the mapping language. \ac{r2rml} is a W3C recommendation that defines how relational tables and \ac{sql} queries can be mapped to \ac{rdf} graphs. The mapping file \texttt{cineexplorer\_mapping.ttl} is serialized in Turtle syntax and processed by the \texttt{r2rml.jar} tool used in the lab sessions. The generated \ac{rdf} output is written to \texttt{output/cineexplorer\_kg.ttl} in Turtle format.

The mapping is executed with the following command from the project root:

\begin{verbatim}
java -jar tools/r2rml/r2rml.jar mapping/mapping.properties
\end{verbatim}

where \texttt{mapping.properties} provides the JDBC connection string, credentials, input mapping file path, and output file path.

The mapping strategy follows the conceptual structure of the ontology rather than directly mirroring the table structure. In particular, the mapping is organized around the main ontology entities (\texttt{CreativeWork}, \texttt{Person}, \texttt{Participation}, and \texttt{Genre}) and around a small number of additional triples maps used to enrich these entities with structural or descriptive information.

The mapping consumes the canonical schema of Section~\ref{sec:infosystem} as published, with no further view layer. Multi-valued attributes (genres on \texttt{Title}, character names on \texttt{Participation}, professions on \texttt{Person}) are already 1NF-flattened into dedicated relations during the \ac{sql} \ac{etl} (Section~\ref{sec:sql}), which means the \ac{r2rml} side does not need any string-splitting tricks: each multi-valued attribute is one logical table and one triples map. This is a deliberate division of responsibilities --- the relational layer handles 1NF, the \ac{r2rml} layer handles class-and-property assignment --- and matches the textbook \ac{erd} reduction discussed in the lectures.

\subsection{IRI Strategy}

A consistent \ac{iri} strategy was defined to ensure that each entity has a stable identifier. \acp{iri} reuse the \ac{imdb} identifiers already present in the source data: \texttt{tconst} values for titles and \texttt{nconst} values for persons. These identifiers are stable, globally recognised, and make it unnecessary to introduce any additional keys into the schema. All \acp{iri} share the project's local namespace, \texttt{http://cineexplorer.local/\allowbreak data/}.

The main \ac{iri} templates used are:

\begin{sloppypar}
\begin{itemize}
    \item \textbf{Creative works} (\texttt{Film}, \texttt{Series}, \texttt{Episode}):\\
\texttt{http://cineexplorer.local/data/title/\{TCONST\}}\\ where \texttt{TCONST} is the \ac{imdb} \texttt{tconst} (e.g., \texttt{tt0111161} for \emph{The Shawshank Redemption}).

    \item \textbf{Persons}:\\
\texttt{http://cineexplorer.local/data/person/\{NCONST\}}\\ where \texttt{NCONST} is the \ac{imdb} \texttt{nconst} (e.g., \texttt{nm0000128} for Russell Crowe).

    \item \textbf{Genres}:\\
\texttt{http://cineexplorer.local/data/genre/\{GENRE\_NAME\}}\\ where \texttt{GENRE\_NAME} is the genre string itself (e.g., \texttt{Drama}). The genre vocabulary is small, bounded, and the name is its own identifier in \ac{imdb}'s data, so no surrogate key is needed.

    \item \textbf{Participation nodes}:\\
\texttt{http://cineexplorer.local/data/\allowbreak participation/\{TCONST\}/\{ORDERING\}}\\ where the composite key \texttt{(TCONST, ORDERING)} matches the primary key of the \ac{sql} \texttt{principal} table (the relational schema uses \ac{imdb}'s source vocabulary; the ontology layer uses \texttt{ce:Participation} for the same concept --- see Section~\ref{sec:infosystem}). \texttt{ORDERING} is the per-title credit position from \texttt{title.principals}, which uniquely identifies one credit row even when the same person holds more than one credit on the same title (the \texttt{ORDERING} differs across credits). \texttt{NCONST} is intentionally not embedded in the \ac{iri}: a participation \ac{iri} keeps its identity even if the credited person is corrected upstream, and the \texttt{playedBy} edge still recovers the link.
\end{itemize}
\end{sloppypar}

The migration tightened this scheme relative to earlier drafts. The previous Participation \ac{iri} template used the triple \texttt{(TITLE\_ID, TALENT\_ID, ORD)}, but with the official \ac{imdb} data it became clear that \texttt{ORD} alone (which \ac{imdb} calls \texttt{ordering}) already disambiguates within a title, and folding \texttt{NCONST} into the \ac{iri} duplicates information already available through \texttt{ce:playedBy}. The two-key form \texttt{\{TCONST\}/\{ORDERING\}} is also exactly the canonical primary key of the \ac{sql} \texttt{principal} relation, which keeps the \ac{iri} strategy aligned with the relational schema.

\subsection{Mapping Structure}

The mapping is organized by entity type and by the role played by each triples map in the generation of the \ac{rdf} graph.

% \paragraph{Creative works.}
\begin{sloppypar}
Three triples maps generate the three subclasses of \texttt{CreativeWork}. The discriminator is the \texttt{TITLE\_TYPE\_ID} column on \texttt{title}: \texttt{Film} corresponds to \texttt{TITLE\_TYPE\_ID = 'movie'}, \texttt{Series} to \texttt{TITLE\_TYPE\_ID IN ('tvSeries',\allowbreak{}'tvMiniSeries', 'tvPilot')}, and \texttt{Episode} to \texttt{TITLE\_TYPE\_ID = 'tvEpisode'} (with the \texttt{title\_episode} join supplying \texttt{partOfSeries}, \texttt{seasonNumber}, and \texttt{episodeNumber}). Each map carries the popularity composite (\texttt{ce:averageRating} from \texttt{AVERAGE\_RATING}, \texttt{ce:numVotes} from \texttt{NUM\_VOTES}) and the descriptive data properties \texttt{ce:workTitle}, \texttt{ce:originalTitle}, \texttt{ce:releaseYear}, \texttt{ce:runtimeMinutes}, and \texttt{ce:isAdult}. The data properties are individually predicated, so a NULL in one column skips only that triple, not the entire subject row, as required by \ac{r2rml}~\S10.3~\cite{r2rml}.
\end{sloppypar}

Title types outside the \texttt{movie} / \texttt{tvSeries}+\texttt{tvMiniSeries}+\texttt{tvPilot} / \texttt{tvEpisode} triad (e.g., \texttt{short}, \texttt{video}, \texttt{tvMovie}, \texttt{tvSpecial}, \texttt{videoGame}) remain in the \texttt{title} relation but receive no \texttt{rdf:type} assertion: in the slice they account for 25 of 5{,}000 titles. They still appear in the \ac{kg} as objects of \texttt{ce:knownFor} and \texttt{ce:participatesIn} (hence inferred to \texttt{ce:CreativeWork} under \ac{rdfs}), which motivates either extending the mapping to cover them or filtering them out at \ac{etl} time.

% \paragraph{Persons.}
All rows from the \texttt{person} table are mapped to \texttt{ce:Person} with their descriptive data properties (\texttt{ce:personName}, \texttt{ce:birthYear}, \texttt{ce:deathYear}). Role-specific subclass membership (\texttt{Actor}, \texttt{Director}, \texttt{Writer}, \texttt{Composer}, \texttt{Editor}) is populated from \emph{two} sources: the per-credit \texttt{principal} table (joined with \texttt{role}) and the career-level \texttt{person\_profession} table (joined with \texttt{profession}). Both feeds use the same person-IRI template, so the resulting subclass-typing triples merge on the same subject. This two-axis approach (Section~\ref{sec:onto-two-axis-discussion}) means a person can be typed as \texttt{ce:Director} either because they have a director credit on some \texttt{principal} row, or because their \ac{imdb} \texttt{primaryProfession} contains \texttt{director}, or both. The \texttt{ce:hasProfession} datatype property is also emitted from \texttt{person\_profession}, so the career axis is queryable as a string attribute as well.

% \paragraph{Participation.}
The \texttt{Participation} triples map selects from the \ac{sql} \texttt{principal} table (joined with \texttt{role} for the role-name string) and generates one \texttt{ce:Participation} instance per credit row. Each instance is linked to its person via \texttt{ce:playedBy} and to its work via \texttt{ce:participatesIn}. The role name is stored as \texttt{ce:participationRole} and the free-text job description as \texttt{ce:participationProperties}. Each instance carries an \texttt{rdfs:label} constructed from the role name and the \ac{imdb} identifiers of the person and title, providing a human-readable identifier for use in \ac{sparql} results and Linked Data browsers.

Character names are drawn from the \ac{sql} \texttt{principal\_character} table and mapped by a dedicated triples map targeting the same participation \ac{iri} template, adding \texttt{ce:characterName} triples without duplicating participation nodes.

% \paragraph{Template-based joins.}
Cross-entity links such as \texttt{ce:partOfSeries} and \texttt{ce:hasGenre} are expressed using matching \texttt{rr:template} patterns on both sides rather than \texttt{rr:RefObjectMap} with explicit join conditions. This is possible because all entities of the same type share a single \ac{iri} template: any map that knows a \texttt{TITLE\_ID} can construct the correct title \ac{iri} without needing a formal join. The advantage is simplicity --- no join overhead and no dependency on the logical table structure. The disadvantage is tight coupling: if the \ac{iri} template for one entity type were to change, every map that references it via template would also need to be updated. An \texttt{rr:RefObjectMap} approach would localise that change to a single parent map.

% \paragraph{Genre links.}
Genres are stored relationally as a multi-valued attribute (\texttt{title\_genre(TCONST, GENRE\_NAME)}), without a separate \texttt{genre} entity table. Two triples maps over \texttt{title\_genre} produce the \ac{rdf} representation: one creates each \texttt{ce:Genre} node from \texttt{DISTINCT GENRE\_NAME} (giving 26 instances in the slice), the other emits the \texttt{ce:hasGenre} edge from each \texttt{Title} to its genre. Both maps share the \ac{iri} template \texttt{.../data/genre/\{GENRE\_NAME\}}, so the two layers naturally merge on the same subjects.

% \paragraph{Known-for links.}
The \texttt{person\_known\_for} table, which records up to four notable works associated with a person (sourced from \ac{imdb}'s \texttt{name.basics.knownForTitles}), is mapped to \texttt{ce:knownFor} links from person \acp{iri} to title \acp{iri}. FK closure is enforced at \ac{etl} time (Section~\ref{sec:sql}), so every emitted \texttt{ce:knownFor} object resolves to a title that is actually present in the \ac{kg}.

% \paragraph{Additional traversal-oriented properties.}
Besides the participation-centered representation, the mapping also materializes several direct properties between persons and creative works, such as \texttt{workedFor} and \texttt{employed}, together with more specific sub-properties such as \texttt{actedIn}, \texttt{directed}, \texttt{wrote}, \texttt{edited}, and \texttt{composedFor}. These properties are redundant from a purely normalization-oriented point of view, since they can be recovered through \texttt{Participation}, but they make graph traversal and query writing more convenient.

The \texttt{workedWith} relation between persons who collaborated on the same work is not materialised in the mapping. It is derivable at query time via shared \texttt{Participation} nodes:

\begin{verbatim}
SELECT DISTINCT ?p1 ?p2 WHERE {
  ?par1 ce:playedBy ?p1 ; ce:participatesIn ?t .
  ?par2 ce:playedBy ?p2 ; ce:participatesIn ?t .
  FILTER(?p1 != ?p2)
}
\end{verbatim}

Materialising this relation via a self-join on \texttt{title\_principal} would generate $O(n^2)$ triples and produce both directions of each pair, which is redundant given the \texttt{owl:SymmetricProperty} declaration on \texttt{ce:workedWith}.

% \paragraph{Language and region metadata.}
Language and region information is derived directly from \texttt{title\_aka} using the ISO codes stored in the \texttt{LANGUAGE\_ID} and \texttt{REGION\_ID} columns. Language and region are attached as ISO-coded string attributes to \texttt{CreativeWork} instances via \texttt{ce:language} and \texttt{ce:region}; the relational \texttt{Language} and \texttt{Region} lookup entities are not promoted to \ac{rdf} classes, since the queries in Section~\ref{sec:sparql} treat them as filter codes rather than entities to navigate.


The generated knowledge graph is serialised to Turtle format and loaded into an Apache Fuseki triplestore for persistent storage and \ac{sparql} querying, as described in Section~\ref{sec:deployment}. Because the ontology targets the \ac{owl}~2 QL profile (Section~\ref{subsec:owl-profile}), the relevant reasoning closures (subclass, sub-property, inverse, domain/range) admit query rewriting and can be applied at query time without materialising all inferred triples beforehand; alternatively, a reasoner can be applied offline to compute closures, or queries can rely on \ac{sparql} property paths for traversal that would otherwise require inference.

% \subsection{Encountered Problems}

% A first difficulty was handling \texttt{NULL} values in the source data. In R2RML, a predicate-object map is skipped when the referenced column is \texttt{NULL}. This is the appropriate behavior for optional properties such as \texttt{releaseYear}, \texttt{runtimeMinutes}, \texttt{birthYear}, and \texttt{deathYear}. As a consequence, the generated RDF simply omits missing values rather than representing their absence explicitly, which is consistent with the open-world assumption.

% A second difficulty concerned datatype assignment. In particular, year-valued properties are represented using \texttt{xsd:gYear}, but the R2RML processor does not infer this datatype automatically from the SQL schema. Explicit \texttt{rr:datatype xsd:gYear} declarations therefore had to be added to the relevant object maps.

% A third difficulty was that some ontology-level relations, such as direct role-specific properties, can be generated in more than one way from the source data. We chose to keep the \texttt{Participation} node as the primary modeling pattern and to treat the direct person--creative work relations as traversal-oriented materializations, which preserves the conceptual clarity of the ontology while improving query usability.